\section{Balances and movements retain their different times}
\label{sec:temporal-statement-network}

A cash-flow statement connects a balance at the start of a period, movements
during that period and a balance at its end. The three quantities have a common
monetary unit but different time meanings. Requiring their complete measurement
types to be identical would reject a valid stock-flow relation; deleting their
dates would lose the reason that the relation is valid. The language therefore
checks the roles of the quantities before solving their numerical equation.

The retained annual example reports cash-flow-statement opening cash of 30,737
at 30 September 2023, a change of \(-794\) during the following fiscal year,
and closing cash of 29,943 at 28 September 2024. The corresponding interim
statement reports an opening amount of 29,943, change of 356 and closing amount
of 30,299 at 28 December 2024. These are source-bound cash-flow captions.
They do not identify the available cash of a particular legal entity or bank
account. Treating their opening and closing captions as one population is an
explicit proposed interpretation in the conditional calculation.

\subsection{A relation can connect different measurement types}

Each temporal measurement declares a value type, a time and a named quantity.
The value type retains issuer, reporting perimeter, reporting basis, monetary
dimension and scale. The time is either a dated snapshot or a half-open period.
The named quantity distinguishes, for example, the cash population reported by
a cash-flow statement from legal-book cash available to meet a payment. The
source mapping preserves the complete original context, including its period
wording and measurement-kind description, alongside this explicit type.

For date-only records, write the flow period as \([s,e)\), with its last
included date immediately before \(e\). The opening snapshot is dated the day
before \(s\), and the closing snapshot the day before \(e\). This is the
convention used by the retained source records, not an intraday accounting rule.
For a declared population \(q\), a stock-flow relation has the form
\begin{equation}
\label{eq:temporal-network-bridge}
 B_q(e-1\text{ day})
 =B_q(s-1\text{ day})+\sum_{j=1}^{m}a_j F_j([s,e)).
\end{equation}
The opening and closing measurements must name the same quantity \(q\).
Every movement must cover the declared period, and all operands must have
compatible monetary value types. The coefficients \(a_j\) retain the source
sign convention. A separately proposed identity specifies why those particular
movements account for the change in the named population.

For the annual cash-flow example, \(s\) is 1 October 2023 and \(e\) is
29 September 2024. For the interim example, they are 29 September and
29 December 2024. Substituting the displayed amounts gives
\begin{equation}
\label{eq:temporal-network-cash-examples}
 30{,}737+(-794)=29{,}943,
 \qquad
 29{,}943+356=30{,}299.
\end{equation}
The negative annual movement reduces the closing balance; the positive interim
movement increases it. These equalities are checked displayed arithmetic.
Whether the same identities hold for the underlying unrounded measurements
requires the explicit identity, population and rounding hypotheses.

The annual closing and interim opening have the same displayed value and date,
but they remain distinct source observations. Their reporting bases differ:
one belongs to the annual issuer report and the other to the condensed interim
statements. The reference neither merges their identities nor imposes equality
between their unrounded amounts merely because the displays agree.

The \path{TemporalMeasurementConstraints/1} reference accepts two relation
forms. A same-time relation requires matching value types and time declarations
for all operands. A stock-flow relation requires the opening, movement and
closing roles just described. Once those checks pass, it lowers the numerical
part of equation~\eqref{eq:temporal-network-bridge} to an opening coefficient
of \(+1\), the supplied movement coefficients and a closing coefficient of
\(-1\). One observation identity still means one unknown throughout the system.
The original temporal definition and the generated scalar constraints are both
retained, so the calculation can be traced back through that lowering.

This does not weaken the earlier same-context profile. It remains an error to
add differently timed quantities through an ordinary same-time relation.
Changing the opening date by one day, swapping opening and closing roles,
using a movement over another period, or changing the endpoint population is
rejected even when the numerical amounts happen to balance. Missing hypotheses
leave a well-formed question unresolved; they do not conceal malformed types.

\subsection{Construct the requested network before asking whether it fits}

The larger retained description has 143 relation records. The adapter first
classifies them by their existing meaning and context, before examining solver
outcomes. Forty snapshot sums, 30 period sums, eight same-period comparisons,
16 same-date endpoint comparisons and ten stock-flow bridges supply 104
candidate equations. Their monetary units and invariant contexts match within
each relation, and their dates fit the applicable temporal form.

Three candidates depend on source concepts omitted from the interim table.
The capital and whole-equity bridges omit the annual concept of equity issued;
the financing-cash relation omits debt issuance. The fact that a row is absent
does not establish a zero movement. Without a separate interpretation for those
omissions, the adapter leaves these three relations outside its supported
selection and preserves the original omission records. It therefore supports
101 conditional equations in this version.

The other 39 source relations also retain their distinct roles: six derived
subtotals, 24 explanatory differences, two endpoint differences, two cash
population comparisons and five cross-period or cross-basis comparisons.
A derived subtotal lacks an independently observed target; an explanatory
difference is not automatically an identity. The population and cross-basis
comparisons require their own interpretations. Classifying these records keeps
their evidence available without silently asserting equations that their
source roles do not justify.

Every selected relation requires named proposed hypotheses for its underlying
identity and population interpretation. Rounding remains separately optional.
The adapter recomputes the complete retained statement description and derives
its coefficients and dates from that description. It accepts neither substitute
source coefficients nor an arbitrary solver result. Original observation digests
provide sharing across equations; equal captions, dates or displayed values do
not merge different observations.

\subsection{The PDF-only request remains unresolved on a reported dash}

Selecting all 101 supported relations uses 236 distinct observations. One of
them, the annual debt-issuance amount in the financing-cash relation, is a
reported dash. Its source interpretation remains nonnumeric. Even with all
proposed identity, population and rounding assumptions supplied, the requested
system is therefore unresolved. The result identifies that observation and
relation. It supplies no witness and does not replace the dash with zero.

There is a useful, different question: are the explicitly selected relations
whose operands have numeric displayed amounts jointly compatible? That selection
omits the annual financing-cash relation in advance, leaving 100 equations and
229 distinct observations. Under the proposed assumptions, this narrower system
has an exact rational witness. Independent checks substitute the witness into
every original source equation and separately check the signed rounding bounds.
The result records the excluded relation and states that the supported network
was not selected in full.

The distinction is substantive. A witness for 100 equations does not resolve
the missing measurement in the 101-equation request. Nor are equations removed
because the solver dislikes them: the narrower request follows the pre-existing
nonnumeric source disposition. If a fully specified selected system is
infeasible, that result concerns its declared assumptions and equations; it is
not permission to drop inconvenient relations until a witness appears.

Both requests preserve all 256 observations and all 143 original dispositions.
Seven observations used only by the omitted annual financing-cash relation
fall outside the narrower numerical system, but remain in the preserved source
description. Dashes, omitted concepts, derived subtotals and explanatory
differences are therefore still available for subsequent interpretation.
The result distinguishes selected relations, supported but unselected relations,
and relations outside this version's capabilities.

\subsection{A second source representation supplies a reported value}
\label{sec:debt-selected-reading}

The unresolved dash concerns the selected PDF reading, not every available
representation of the filing. A later inspection of the retained annual XBRL
instance finds a numerical fact for the same caption. Occurrence
\code{f-325}, under \code{ProceedsFromIssuanceOfLongTermDebt}, reports zero
USD for 1 October 2023 through 28 September 2024. It identifies Apple's CIK,
has no additional segment or scenario dimensions, and declares decimal
accuracy \code{-6}. The filing-specific label is \emph{Proceeds from issuance of
term debt, net}, matching the annual cash-flow statement. Its presentation
relationship places it in financing activities, and the calculation
relationship gives it a positive contribution to the financing total
\citep[retained instance, occurrence f-325, context c-1; associated label,
presentation and calculation relationships]{apple2024xbrlinstance}.

This is a separate source-linked reading. The PDF still contains a dash;
its original observation and the earlier unresolved calculation remain intact.
The selected interpretation retains both representations and their exact
identities. A dash elsewhere, an absent interim row, a nil fact and a
reported numerical zero do not become interchangeable. The XBRL value's
declared accuracy also matters: a reported zero at million-dollar accuracy
does not establish an exact underlying economic zero
\citep[Component Fact, value and decimals properties]{xbrloim}.

The neighboring years and repayment row provide useful checks on the match.
The same issuance concept reports USD~5,228 million and USD~5,465 million in
the two comparative years, agreeing with the visible columns. The FY2024
repayment fact reports USD~9,958 million as a payment; its negative
calculation weight agrees with the parenthesized cash-flow amount. All
thirty retained occurrences of the eight financing concepts, including
duplicates, stay distinct. The seven signed FY2024 components then sum to
\(-5{,}441-15{,}234-94{,}949+0-9{,}958+3{,}960-361=-121{,}983\)
million dollars, matching the displayed financing total
\citep[p.~33]{apple2024}. The earlier no-contribution convention is therefore
no longer the only available basis for this displayed calculation: the
selected numerical reading has its own inspected source support.

The term-debt comparison remains a different question. Closing current and
noncurrent carrying amounts sum to USD~96,662 million, compared with
USD~105,103 million at the opening date. Subtracting the selected reported
issuance and the signed repayment gives
\(96{,}662-105{,}103-0-(-9{,}958)=1{,}517\) million dollars.
This is an unattributed difference between carrying-value movement and the
selected cash flows. It supplies no allocation to foreign exchange, hedges,
fees or another account and creates no balancing journal.

The selected interpretation can also be carried into the whole conditional
measurement system. It changes one unknown interval while preserving all
236 selected observations and all 101 equations. The complete parent still
retains the PDF dash, all 256 source observations and all 143 relation
dispositions. The rounding rule remains the earlier, explicitly proposed
nearest-million rule with half ties away from zero. Under that assumption,
the latent issuance amount for a reported zero lies in the open interval
\((-1/2,1/2)\) million dollars. This interval comes from the chosen rounding
rule; the source accuracy alone is not used to impose its endpoints or an
additional nonnegativity condition.

The full conjunction admits an exact rational witness. A separate checker
substitutes every value into each original equation and interval. A second
complete witness fixes the issuance measurement at USD~0.25 million and
still satisfies the same original constraints. This is a useful discriminator:
the numerical reading resolves the selected source uncertainty without
pretending that the latent amount is known exactly. No financing equation is
dropped to obtain either witness.

Missing interpretation or rounding premises remain unresolved. A changed
issuer, period, source identity, original dash, equation set or rounding rule
is refused; an exhausted computation budget remains unknown. The interface
retains the original PDF-only result separately. These checks establish a
conditional measurement composition, not an actual sequence of company events.
Legal books, instrument populations, source availability and the remaining
identity and population assumptions still require their own evidence.

The exact source selection and signed comparisons are retained under
\path{data/accounting-master/debt-dash-source-20261027/}; the full network,
independent witness checks and refusal cases are under
\path{data/accounting-master/debt-reading-network-20261027/}.

\subsection{What the temporal network contributes to reconstruction}

The selected network connects balance-sheet, income, comprehensive-income,
cash-flow and equity measurements using shared identities and explicit times.
It extends the earlier isolated balance equations without inventing a sequence
of company transactions. The recorded checks include altered dates, periods,
source coefficients, quantity interpretations, hypotheses and completeness
claims. They also compare the exact terms and intervals with the preceding
balance and cash-bridge references. A zero computational budget returns unknown,
and a valid alternative rational witness is allowed on replay.

The implementation and complete retained examples are under
\path{data/accounting-master/}:
\begin{quote}\small
\path{temporal-measurement-contexts-20261024}\\
\path{statement-network-contexts-20261024}
\end{quote}
They establish checked conditional measurement behavior. Their interpretation
still depends on proposed identities and populations; feasibility does not
promote those premises to issuer facts or accounting authority.

Section~\ref{sec:conditional-history-families} develops a finite conditional
connection to authored states and event sequences. Actual historical
reconstruction still requires the corresponding source and method admission. Several histories can produce the same balances
and movements while leaving different contracts, tax attributes, inventory
cohorts or liquidity constraints at the end. Those differences matter for the
simulator's next event. A history-family representation must preserve them,
state its search scope and retain uncomputed or missing cases. Until that
connection is justified, the retained company has no executable legal-book
opening and its actual history remains unknown.
